% SEGMENT OLP-0067-S01
% Part: first-order-logic
% Chapter: proof-systems
% Section: tableaux

\documentclass[../../../include/open-logic-section]{subfiles}

\begin{document}

\iftag{FOL}
      {\olfileid{fol}{prf}{tab}}
      {\olfileid{pl}{prf}{tab}}

\olsection{\usetoken{P}{tableau}}

பல !!{derivation} முறைமைகள் !!{sentence}s இன் ஒழுங்கமைப்புகளைக் கொண்டு
செயல்படுகின்றன; !!{tableau}s, !!{signed formula}s ஐக் கொண்டு
செயல்படுகின்றன. !!^a{signed formula} என்பது ஒரு மெய்மதிப்புக் குறியையும்
($\True$ அல்லது $\False$) !!a{sentence}வையும் கொண்ட சோடி:
\[
\sFmla{\True}{!A} \text{ அல்லது } \sFmla{\False}{!A}.
\]
!!^a{tableau} என்பது கீழ்நோக்கிக் கிளைக்கும் ஒரு மரத்தில் ஒழுங்கமைக்கப்பட்ட
!!{signed formula}s ஐக் கொண்டது. அது சில \emph{எடுகோள்களுடன்} தொடங்கி,
அவற்றுக்கு மேலுள்ள !!{signed formula}s ஒன்றிற்கு அனுமான விதி ஒன்றைப்
பயன்படுத்துவதால் கிடைக்கும் !!{signed formula}s உடன் தொடர்கிறது. ஒவ்வொரு
விதியும் ஒரு கிளையின் முடிவில் ஒன்று அல்லது அதற்கு மேற்பட்ட !!{signed
  formula}s ஐச் சேர்க்க அனுமதிக்கும்; அல்லது இரண்டு !!{signed formula}s ஐ
அருகருகே சேர்க்க அனுமதிக்கும். பின்னைய நேர்வில், ஒரு கிளை இரண்டாகப்
பிரிந்து, சேர்க்கப்பட்ட இரண்டு !!{signed formula}s அந்த இரு கிளைகளின்
முடிவுகளாகின்றன.

% SEGMENT OLP-0067-S02
ஒரு கூட்டு !!{signed formula}வுக்குப் பயன்படுத்தப்படும் விதி, அதன் உடனடித்
துணை-!!{formula}s ஆன !!{signed formula}s ஐச் சேர்க்கிறது. அவை சோடிகளாக
வருகின்றன: இரு குறிகளில் ஒவ்வொன்றுக்கும் ஒரு விதி. எடுத்துக்காட்டாக,
$\TRule{\True}{\land}$ விதி $\sFmla{\True}{!A \land !B}$ க்குப்
பொருந்துகிறது; $\sFmla{\True}{!A}$ மற்றும்~$\sFmla{\True}{!B}$ என்ற இரு
!!{signed formula}s ஐயும், $\sFmla{\True}{!A \land !B}$ ஐக் கொண்ட எந்தக்
கிளையின் முடிவிலும் சேர்க்க அது அனுமதிக்கிறது. $\TRule{\False}{!A \land !B}$
விதி $\sFmla{\False}{!A}$ மற்றும் $\sFmla{\False}{!B}$ ஐ அருகருகே
சேர்த்துக் கிளையைப் பிரிக்க அனுமதிக்கிறது. ஓர் !!^{tableau}வின் ஒவ்வொரு
கிளையிலும் $\sFmla{\True}{!A}$ மற்றும் $\sFmla{\False}{!A}$ என்ற
பொருந்தும் !!{signed formula}s சோடி இருந்தால் அந்த அட்டவணை மரம் மூடியது.

% SEGMENT OLP-0067-S03
!!{tableau}s ஐ அடிப்படையாகக் கொண்ட $\Proves$ உறவு பின்வருமாறு
வரையறுக்கப்படுகிறது: எடுகோள்களின்
\[
\{\sFmla{\False}{!A}, \sFmla{\True}{!B_1}, \dots, \sFmla{\True}{!B_n}\}
\]
கணத்திற்காக ஒரு மூடிய !!{tableau} இருக்குமாறும், $\Gamma_0 = \{!B_1,
\dots, !B_n\} \subseteq \Gamma$ ஆகுமாறும் ஏதோ ஒரு முடிவுறு கணம்
இருந்தால், அப்போதும் அப்போது மட்டுமே, $\Gamma \Proves !A$.
எடுத்துக்காட்டாக, $\Proves (!A \land !B) \lif !A$ என்பதைக் காட்டும் ஒரு
மூடிய !!{tableau} இதோ:
\begin{oltableau}
  [\sFmla{\False}{(\formula{A} \land \formula{B}) \lif \formula{A}}, just = \TAss
    [\sFmla{\True}{\formula{A} \land \formula{B}}, just = {\TRule{\False}{\lif}[1]}
      [\sFmla{\False}{\formula{A}}, just = {\TRule{\False}{\lif}[1]}
        [\sFmla{\True}{\formula{A}}, just = {\TRule{\True}{\lif}[2]}
          [\sFmla{\True}{\formula{B}}, just = {\TRule{\True}{\lif}[2]}, close]
        ]
      ]
    ]
  ]
\end{oltableau}

% SEGMENT OLP-0067-S04
ஏதோ சில $!B_i \in \Gamma$ க்காக, எடுகோள்களின்
\[
\{\sFmla{\True}{!B_1}, \dots, \sFmla{\True}{!B_n}\}
\]
கணத்திற்கு ஒரு மூடிய !!{tableau} இருந்தால், கணம் $\Gamma$ !!{tableau}
சார்ந்த கணியத்தில் முரணுடையது.

% SEGMENT OLP-0067-S05
!!^{tableau}s ஐ 1950களில் எவர்ட் பெத் மற்றும் யாக்கோ ஹிண்டிக்கா தனித்தனியே
உருவாக்கினர்; ரேமண்ட் ஸ்முல்லியன் அவற்றை எளிமைப்படுத்திப் பரவலாக்கினார்.
!!a{tableau}வை அமைப்பது மிகவும் முறையான செயல்முறை என்பதால், அவற்றைப்
பயன்படுத்துவது மிக எளிது. !!{tableau}s இன் முறையான இயல்பினால், அவற்றைக்
கணினியிலும் செயல்படுத்த முடியும். இருப்பினும், !!a{tableau}வை வாசிப்பது
பெரும்பாலும் கடினம்; நிறுவல்களுடன் அவற்றுக்குள்ள தொடர்பையும் சிலவேளைகளில்
எளிதாகக் காண முடியாது. இந்த அணுகுமுறை மிகவும் பொதுவானது; பல வெவ்வேறு
தருக்கவியல்களுக்கு !!{tableau} முறைமைகள் உள்ளன. கொடுக்கப்பட்ட !!{sentence}s
கணத்தை நிறைவுறுத்தும் !!{structure}s ஐக் கண்டறியவும் !!^{tableau}s
உதவுகின்றன: அந்தக் கணம் நிறைவுறத்தக்கது எனில், அதற்கு ஒரு மூடிய
!!{tableau} இருக்காது; அதாவது, எந்த !!{tableau}விலும் ஒரு திறந்த கிளை
இருக்கும். அந்தக் கிளையில் பயன்படுத்தக்கூடிய ஒவ்வொரு விதியும் அதில்
பயன்படுத்தப்பட்டிருந்தால், நிறைவுறுத்தும் !!{structure}வை அந்தத் திறந்த
கிளையிலிருந்து ``வாசித்தெடுக்கலாம்.'' தொடரணிக் கணியத்துடனும் மிக நெருங்கிய
தொடர்பு உண்டு: சாரத்தில், ஒரு மூடிய !!{tableau} என்பது தலைகீழாக எழுதப்பட்ட,
தொடரணிக் கணியத்தின் சுருக்கப்பட்ட !!{derivation} ஆகும்.

\end{document}
